% \iffalse meta-comment
%
% Copyright (c) 2026- Jesse Straat
% 
% This work may be distributed and/or modified under the
% conditions of the LaTeX Project Public License, either version 1.3
% of this license or (at your option) any later version.
% The latest version of this license is in
%   https://www.latex-project.org/lppl.txt
% and version 1.3c or later is part of all distributions of LaTeX
% version 2008 or later.
%
% This work has the LPPL maintenance status `maintained'.
% 
% The Current Maintainer of this work is Jesse Straat.
%
% This work consists of the main file collabcomment.dtx
% and the derived files
%      collabcomment.sty, collabcomment.pdf, collabcomment.ins
%
% Unpacking:
% (a) If collabcomment.ins is present:
%		pdflatex collabcomment.ins
% (b) If collabcomment.ins is absent:
%		pdflatex collabcomment.dtx
%
% Documentation:
% 	    pdflatex collabcomment.dtx
%       makeindex -s gind.ist -o collabcomment.ind collabcomment.idx
%       makeindex -s gglo.ist -o collabcomment.gls collabcomment.glo
%       pdflatex collabcomment.dtx
%
%<*ignore>
\iffalse
%</ignore>
%
%<*readme>
# `collabcomment`

PLACEHOLDER README

# Installation
Run the following in the command line:
1\. If collabcomment.ins is present:
    `pdflatex collabcomment.ins`
2\. If collabcomment.ins is absent:
    `pdflatex collabcomment.dtx`

Move `.sty` and `.cls` files to a folder that TeX can find.

# Documentation
Run the following in the command line:
```
pdflatex collabcomment.dtx
makeindex -s gind.ist -o collabcomment.ind collabcomment.idx
makeindex -s gglo.ist -o collabcomment.gls collabcomment.glo
pdflatex collabcomment.dtx
```
This automatically unpacks the package, as well.


If you are looking for comments or documentation, you won't find
any in the source. These packages make use of literate programming, which means
that the documentation is given in the form of a pdf file. You can
probably find collabcomment.pdf in this folder, by running "texdoc --view collabcomment"
in the command line, or Googling
"[your distribution] how to find package documentation".
If you're using MiKTeX, open MiKTeX console, go to
"Documentation" and look for "collabcomment".

---

&copy; 2026- Jesse Straat

License: [LPPL1.3c](https://www.latex-project.org/lppl.txt)

[GitHub Repository](https://github.com/JesseStraat/collabcomment)

%</readme>
%
%<*ignore>
\fi
\def\plainTeXstring{plain}
\ifx\fmtname\plainTeXstring\else
  \expandafter\begingroup
\fi
%</ignore>
%
%<*install>
\input docstrip.tex
\Msg{************************************************************************}
\Msg{* Installation}
\Msg{* Package: collabcomment 2026/06/02 v1.0}
\Msg{************************************************************************}

\keepsilent
\askforoverwritefalse

\preamble

Copyright (c) 2026- Jesse Straat

This work may be distributed and/or modified under the
conditions of the LaTeX Project Public License, either version 1.3
of this license or (at your option) any later version.
The latest version of this license is in
  https://www.latex-project.org/lppl.txt
and version 1.3c or later is part of all distributions of LaTeX
version 2008 or later.

This work has the LPPL maintenance status `maintained'.

The current maintainer of this work is Jesse Straat.

This work consists of the main file collabcomment.dtx
and the derived files
    collabcomment.sty, collabcomment.pdf, collabcomment.ins

If you are looking for comments or documentation, you won't find
any in the source. These packages make use of literate programming, which means
that the documentation is given in the form of a pdf file. You can
probably find PLACEHOLDER.pdf in this folder, by running "texdoc --view collabcomment"
in the command line, or Googling
"[your distribution] how to find package documentation".
If you're using MiKTeX, open MiKTeX console, go to
"Documentation" and look for "collabcomment".

\endpreamble

\usedir{tex/latex/collabcomment}
\generate{
    \file{collabcomment.sty}{\from{\jobname.dtx}{main}}
}

\obeyspaces
\Msg{************************************************************************}
\Msg{*}
\Msg{* To finish the installation you have to move the following}
\Msg{* files into a directory searched by TeX:}
\Msg{*}
\Msg{*    collabcomment.sty}
\Msg{*}
\Msg{* To produce the documentation run the file `collabcomment.dtx'}
\Msg{* through LaTeX.}
\Msg{*}
\Msg{* Happy TeXing!}
\Msg{*}
\Msg{************************************************************************}

%</install>
%<install>\endbatchfile
%
%<*ignore>
\usedir{source/latex/collabcomment}
\generate{
    \file{\jobname.ins}{\from{\jobname.dtx}{install}}
}

\nopreamble\nopostamble
\usedir{doc/latex/collabcomment}
\generate{
    \file{README.md}{\from{\jobname.dtx}{readme}}
}

\ifx\fmtname\plainTeXstring
  \expandafter\endbatchfile
\else
  \expandafter\endgroup
\fi
%</ignore>

%<*driver>
\NeedsTeXFormat{LaTeX2e}[2020/10/01]
\documentclass{ltxdoc}
\usepackage{\jobname}
\usepackage[english]{babel}
\usepackage[final,
hyperindex=false,
colorlinks,
]{hyperref}
\EnableCrossrefs
\CodelineIndex
\RecordChanges
\begin{document}
  \DocInput{\jobname.dtx}
\end{document}
%</driver>
% \fi
%
% \GetFileInfo{\jobname.sty}
%
% \DoNotIndex{\begin,\end}
%
% \title{^^A
%     collabcomment\thanks{^^A
%         Version \fileversion, last revised \filedate.^^A
%     }^^A
% }^^A
% \author{^^A
% 	Jesse Straat^^A
% }
% \date{\filedate}
%
% \maketitle
%
% \begin{abstract}
%     \noindent Abstract
% \end{abstract}
%
% \tableofcontents
%
% \changes{v1.0}{2026/06/02}{First release}
%
%
%
%\StopEventually{^^A
%  \clearpage
%  \PrintChanges
%  \PrintIndex
%}
%
%
%
%
% \section{Documentation}
% Bla
%^^A
% \subsection{Usage}
% \createauthor{test}{blue}
% \listauthors
% Bla
% \todo[test][margin]{Test!!}
% \todo[fake]{Test?}
% \begin{equation}
%     E = \todo[test]{\gamma}mc^2
% \end{equation}
% \todo*[test]{This comment is done}
% \todolist
% 
%
%
%
% \clearpage
% \section{Source code}
% This section contains the source code of the library.
% It contains some explaining descriptions, allowing for
% \href{https://https://en.wikipedia.org/wiki/Literate_programming}{literate programming}.
% This is where the documentation for ``regular people'' ends.
%^^A
%^^A
%^^A
%^^A
% \setcounter{CodelineNo}{0}
%    \begin{macrocode}
%<*main>
%<@@=collabcomment>
\ProvidesPackage{collabcomment}[2026/06/02 v1.0 Collaborative comments]
\RequirePackage{xcolor}
\RequirePackage{xparse}
\ExplSyntaxOn

\prop_new:N \@@_prop_authorcolours
\prop_set_from_keyval:Nn \@@_prop_author_colours {
    default = red
}
\prop_new:N \@@_prop_comment_list
\prop_set_eq:NN \@@_prop_comment_list \c_empty_prop
\prop_new:N \@@_prop_todo_list
\prop_set_eq:NN \@@_prop_todo_list \c_empty_prop

\keys_define:nn {position}{
    margin .bool_set:N = \l_tmpa_bool,
}

\NewDocumentCommand{\comment}{ O{default} O{} +m }{
    % o: Commenter name
    % o: Key-value options
    % m: Comment content
    \@@_parse_keyvals:nn {#2} {\@@_make_comment:nn {#1} {#3}}
    \@@_prop_add_to:Nnn \@@_prop_comment_list {#1} {#3}
}

\NewDocumentCommand{\todo}{ s O{default} O{} +m }{
    % s: Makes the todo "done"
    % o: Commenter name
    % o: Key-value options
    % m: Comment content
    \IfBooleanF{#1}{
        % Todo isn't done yet: display it
        \@@_parse_keyvals:nn {#3} {\@@_make_comment:nn{#2}{\textbf{Todo:}~#4}}
    }
    \@@_prop_add_to:Nnn \@@_prop_todo_list {#2} {
        \IfBooleanT{#1}{\color{lightgray}}
        #4
    }
}

\NewDocumentCommand{\createauthor}{ m m }{
    % m: Author name
    % m: Author colour
    \prop_put:Nnn \@@_prop_author_colours {#1} {#2}
}

\NewDocumentCommand{\listauthors}{}{
    List~of~authors:
    \begin{itemize}
        \prop_map_inline:Nn \@@_prop_author_colours {
            \item
            \begingroup
            \color{##2}
            ##1:~##2\par
            \endgroup
        }
    \end{itemize}
}

\NewDocumentCommand{\commentlist}{}{
    \@@_display_list:N \@@_prop_comment_list
}

\NewDocumentCommand{\todolist}{}{
    \@@_display_list:N \@@_prop_todo_list
}

\cs_new:Nn \@@_make_comment:nn {
    % #1: Commenter name
    % #2: Content
    \begingroup
    \color{\@@_get_author_colour:n{#1}}#2
    \endgroup
}

\cs_new:Nn \@@_get_author_colour:n {
    % #1: Author name
    \prop_if_in:NnTF \@@_prop_author_colours {#1} {
        \prop_item:Nn \@@_prop_author_colours {#1}
    }{
        \prop_item:Nn \@@_prop_author_colours {default}
    }
}

\cs_new:Nn \@@_prop_add_to:Nnn {
    % #1: Comment/todo list name
    % #2: Author
    % #3: Content
    \tl_set:Ne \l_tmpa_tl {\prop_item:Nn #1 {#2}}
    \tl_put_right:Nn \l_tmpa_tl {\item #3}
    \prop_put:NnV #1 {#2} \l_tmpa_tl
}

\cs_new:Nn \@@_display_list:N {
    % #1: comment/todo list name
    \prop_map_inline:Nn #1 {
        \tl_set:Nn \l_tmpa_tl {##2}
        \tl_if_empty:NF \l_tmpa_tl {
            \par
            {\Large \color{\@@_get_author_colour:n {##1}}##1:}
            \begin{itemize}
               ##2
            \end{itemize}
        }
    }
}

\cs_new:Nn \@@_parse_keyvals:nn {
    % #1: Keyval list
    % #2: Input text
    \group_begin:
        \keys_set:nn {position} {#1}
        \bool_if:NTF {\l_tmpa_bool} {
            % Margin option
            \marginpar{#2}
        }{
            % No option
            #2
        }
    \group_end:
}

%    \end{macrocode}
%    \begin{macrocode}
\ExplSyntaxOff
%</main>
%    \end{macrocode}
% \Finale